\input{../tex-inputs/latex-leseansicht-vorspann}

\section[Arthur Schnitzler an Max Schwarzkopf, 13. 10. 1905]{L04017 Arthur Schnitzler an Max Schwarzkopf, 13. 10. 1905}
\nopagebreak\mylabel{L04017v}
\rehead{ }\normalsize\beginnumbering\briefempfaengerindex{Schwarzkopf, Max@\textsc{Schwarzkopf, Max}!zzzSchnitzler, Arthur@\emph{von Arthur Schnitzler}!1905-10-131@{13. 10. 1905}|(be}
\toendnotes[C]{\smallbreak\pagebreak[2]}
\correspDesc{Versand  durch Arthur Schnitzler am 13. 10. 1905 in Wien
\newline{}Erhalt  durch Max Schwarzkopf am 13. 10. 1905 in Wien}\toendnotes[C]{\smallbreak}
\Standort{CUL, Schnitzler, B 96.}
\physDesc{Postkarte, 232 Zeichen
\newline{}Handschrift: schwarze Tinte, deutsche Kurrent
\newline{}Versand: 1) Rohrpost  2) Stempel: »\nobreak{}\oindex{XVIII., Währing@\textbf{XVIII., Währing}, \emph{Verwaltungsgebiet}|pwk}Wien 18/1 110, 13. X. 05, 8.40V\nobreak{}«.  3) Stempel: »\nobreak{}\oindex{I., Innere Stadt@\textbf{I., Innere Stadt}, \emph{Verwaltungsgebiet}|pwk}Wien 1/1, 13. X. 05, 9–V\nobreak{}«. }\toendnotes[C]{\smallbreak}\pstart{}{\pb}Dr. Max
                  Schwarzkopf\pend{}\pstart{}Wien I\oindex{I., Innere Stadt@\textbf{I., Innere Stadt}, \emph{Verwaltungsgebiet}|pw}\pend{}\pstart{}Tiefer Graben 23\oindex{Wien@\textbf{Wien}!I., Innere Stadt@\textbf{I., Innere Stadt}!Tiefer Graben 23@\textbf{Tiefer Graben 23}, \emph{Wohngebäude}|pw}.\pend{}{\bigskip}\vspace{1em}
\pstart
           \noindent{}{\pb}{\pb}lieber Herr Doktor,
               Sie erhalten die Sitze zur heut\textcolor{gray}{g}{ }Vorstellung\eventindex{Burgtheater@\textbf{Burgtheater}!2. Aufführung von Zwischenspiel, 13.10.1905@2. Aufführung von Zwischenspiel, 13.10.1905|pwv} durch die Intendanz\orgindex{Burgtheater@Burgtheater|pwv} zugeſandt. Den 2ten Sitz bitte nach Belieben zu verwenden.\pend
           
\pstart
           Herzliche Grüße an Sie u Guſtav\pwindex{Schwarzkopf, Gustav 7.\,11.\,1853 Wien – 13.\,11.\,1939 ebd.@\textsc{Schwarzkopf, Gustav} (7.\,11.\,1853 Wien – 13.\,11.\,1939 ebd.), \emph{Schriftsteller}|pw}.\pend
           
\pstart
           Ihr{\\[\baselineskip]}\spacefill\mbox{ArthSch}\pend
           \leftskip=0em{}
\pstart
           13. X. 905\pend
           \selectlanguage{ngerman}\endnumbering\briefempfaengerindex{Schwarzkopf, Max@\textsc{Schwarzkopf, Max}!zzzSchnitzler, Arthur@\emph{von Arthur Schnitzler}!1905-10-131@{13. 10. 1905}|)be}\mylabel{L04017h}
\begin{anhang}
\end{anhang}\newcommand{\dateiname}{L04017}\newcommand{\titel}{Arthur Schnitzler an Max Schwarzkopf, 13. 10. 1905}\newcommand{\editorInnen}{Herausgegeben von Jahnke, SelmaMüller, Martin Anton}\input{../tex-inputs/latex-leseansicht-abspann}
